\section{The historical likelihood referent}
\label{app:referents}

The 2016 discussion of Shomorony et al.\ asks whether bridging conditions
guarantee that the \emph{maximum-likelihood sequence} is the true sequence,
citing Medvedev and Brudno (2009)
\parencite{shomorony2016,medvedev2009}. The sentence does not identify a
particular objective, candidate class, strand convention, or tie rule.
These distinctions affect the mathematical statement. We therefore
give separate, explicitly scoped results rather than advertising a single
unqualified resolution of the historical question.

\begin{table}[htbp]
\centering\small
\begin{tabular}{@{}p{0.39\textwidth}p{0.45\textwidth}@{}}
\toprule
Possible referent & Mathematical object being optimized \\
\midrule
Exact multinomial likelihood & A circular genome with candidate-intrinsic length \\
Fixed-$N$ product-binomial surrogate & A genome scored using a fixed external length \\
MB09 bidirected-flow optimizer & A possibly disconnected flow, not necessarily one sequence \\
General maximum-likelihood principle & Candidate set and objective left unspecified \\
\bottomrule
\end{tabular}
\caption{Four distinct interpretations of the cited likelihood formulation.
The operational assembler uses a separable surrogate; the exact
read-count law is a separate objective.}
\label{tab:referents}
\end{table}

\begin{table}[htbp]
\centering\small
\begin{tabular}{@{}p{0.34\textwidth}p{0.50\textwidth}@{}}
\toprule
Interpretive axis & What is established \\
\midrule
Read type & Oriented windows and reverse-complement molecule classes differ \\
Genome length & External binomial reference size does not fix candidate length \\
Maximality & Truth as a maximizer is weaker than uniqueness \\
Genome equivalence & Rotations necessarily tie; molecule types also force reverse-complement ties \\
Flow versus spelling & Balance does not imply a single spelled genome \\
\bottomrule
\end{tabular}
\caption{Axes that must be fixed when interpreting a likelihood claim.}
\label{tab:axes}
\end{table}

The original historical read-string bridging rule has been checked against
the author-hosted supplement. The publisher-hosted supplementary archive
has not been independently retrieved, and the exact relation between
Medvedev--Brudno's $k$-molecule vocabulary and its $4^k$ index count
is unresolved. This limits attribution to the original question, not
the validity of a theorem proved for a named formal model.

The detailed source analysis, quotations, and primary-source citations
remain in \path{docs/source-notes/} and the preserved manuscript source
\path{paper/technical-notes/}. The model
restrictions that matter for each result are stated in the relevant
theorem, rather than being left for this appendix to supply.
